\section{Introduction}
A causal experiment acquires reports through physical preparations and interventions. A prediction must be fixed before its future audit, and failed or rejected observations still consume resources. The central problem of General Theta Foundations I is to connect the predictive information actually acquired by such an experiment to the performance of an implementable finite-resource observer. This requires more than quantizing a prescribed probability law: one must establish an all-strategy converse, construct one program that works in every admitted model, and count its persistent memory without supplying an unobserved state or a free clock.

This paper develops two complementary routes. The \emph{geometric route} identifies an interpretable resource curve from unnormalized task comparisons on actual histories. The \emph{primitive optimization route} starts with the physical instruments and constructs lower and upper approximations to their constrained optimal risk, without assuming that a geometric comparison profile is exact. The former applies to broadly prepared causal kernels with stated certificates. The latter admits a compact continuum of unknown fixed models on a finite-dimensional, finite command/report interface. Neither route silently supplies the assumptions of the other.

\subsection*{Primitive synthesis with an unknown operating model}
Consider a compact continuous family of time-homogeneous classical or quantum instruments with common legal commands and productive reports, bounded nonnegative terminal loss, and call costs between two fixed positive bounds. The initial physical state and the fixed model are unknown. A common dark subspace in any one model prevents proper all-initial-state operation. Conversely, its absence throughout the compact family supplies a single full-support randomized command row with a uniform finite cost, by a quantitative finite-depth test. This uses established quantum reachability geometry \cite{Ying14pre}; compact continuity and the explicit spectral margin make its model-uniform cost certificate available.

A finite set of candidate models is not enough by itself. Even when every primitive entry is affine in a parameter on an interval, testing only its endpoints can give a zero-risk certificate for a problem whose true robust risk is $1/2$; Example~\ref{ex:vertex-risk} proves this for next-bit prediction with a fixed two-value decoder. A robust certificate needs both the finite test and a bound on what was omitted.

Theorem~\ref{thm:compact-risk} supplies that bound. A depth-$T$ common realization plan optimized on a certified model cover has value $q_{F,C,T}(N)$. Its constraints charge the continuing mass for at least one more physical call. Completing the plan by a derived fallback and mixing once at initialization yields a finite-label observer at exactly the original expected-cost budget $N$, simultaneously for every model and initial state. The enclosure is
\[
 q_{F,C,T}(N)\le R_{\Theta,C}(N)
 \le q_{F,C,T}(N)+z_{F,T}
       +\frac{Le_{F,T}}{N-\kappa+e_{F,T}},
 \qquad N>\kappa.
\]
Here $\kappa$ is the actual uniform fallback cost, $z_{F,T}$ is a bounded-task error, and $e_{F,T}$ is the completion and model-induced cost excess. Their explicit formulas are given in \eqref{eq:compact-budget-error}--\eqref{eq:compact-risk-error}. The denominator records the real cost slack needed to restore the budget; it is not a formal rescaling of the loss.

The key comparison, Lemma~\ref{lem:occupation-perturbation}, integrates one-step instrument discrepancies against the tested controller's actual pre-call occupation. At cost $N$, bounded terminal loss incurs at most $L(N/c_*)\epsilon_I+\epsilon_F$ model error, rather than a generic $T\epsilon_I$ error. Cost still contains a remaining-horizon factor, and Example~\ref{ex:cost-horizon} realizes it exactly. Joint refinement with $T\epsilon_I\to0$ therefore gives convergent same-budget risk bounds. A fixed interior slack gives an $O(\sqrt\epsilon)$ upper continuity modulus; its optimal exponent is not asserted. The scalar boundary example proves that continuity can fail at the least feasible budget even when both perturbed instances remain feasible.

The preliminary cost enclosure in Theorem~\ref{thm:compact-cost} converges from both sides to the least feasible uniform cost and constructs a cheaper finite fallback whenever the requested budget is strictly larger. Thus the risk result covers \emph{every interior feasible budget}, not only those above an arbitrary exploratory policy's cost. Finite decision nets keep the same output-width constraint. For compact semialgebraic parameter and decision sets with polynomial primitive entries, Corollary~\ref{cor:compact-algebraic} gives a terminating exact-description synthesis procedure and proves computability of the least feasible cost. This is not a polynomial-time claim or a finite-memory attainment theorem at the boundary.

These conclusions concern the constrained optimum, not only the evaluation of a supplied policy. Their construction does not use a balancing group, a physical reset, a model-identification oracle, or a normalized-filter contraction. The returned observer sees only its actual reports and counted labels. A model fixed for the whole run is essential: allowing an adversary to change it after every call defines a different experiment, as Example~\ref{ex:compact-limits} shows.

\subsection*{Actual-history transfer and acquired geometry}
For a genuinely prepared comparison law, let $f(x)=\E_0[Z\mid x]$, $w(x)=\E_0[\|Z-f(x)\|^2\mid x]$, and let $\nu$ be its report law. The marked profile
\[
 \Q_m(b)=\min_{|C|\le m,\ 0\le h\le b,\ \int h\,d\nu=1}
            \int[w+\dist(f,C)^2]h\,d\nu
\]
is population trimmed quantization with retained mass $1/b$ and a residual-loss mark \cite{CGM97,GG99}. Neither this functional nor ordinary rejection sampling is claimed as a new principle. Its transfer to controlled experiments requires a commandwise full-current-report comparison in one actual operating model, an all-program effort lower, actual-occupation defect control, and a common all-model implementation with counted live and export roles.

Lemma~\ref{lem:comparison} constructs the accepted comparison measure by replacing only the current row inside the competitor's actual history law. It never runs that competitor under a fictitious reference experiment where it could fail to terminate. Theorem~\ref{thm:transfer} then gives the converse and common conditional-hazard upper without a minimax exchange. Relative task order preserves exact zero-task directions and all resolution scales; absolute one-call error can instead grow linearly with the number of acquired samples. The task is fixed in this comparison, not changed between its two directions.

Exact finite-width defects in Theorems~\ref{thm:width-defect} and~\ref{thm:quadratic-defect} avoid charging a finite decoder for an unrestricted family of readouts. They distinguish ordinary squared prediction from the larger affine-square matrix family; equality of their zero-defect orders is not an unqualified Blackwell comparison. The width amplification $\min(m,k)/(5k)$ is sharp in its explicit dimension-$k$ family. Theorem~\ref{thm:width-transfer} extends the transfer to bounded decision-continuous losses, including Bregman prediction. Theorem~\ref{thm:zero-spectrum} determines the exact finite-width zero-risk thresholds from the largest atomic masses of the zero-residual predictive measure, correcting the temptation to use one favorable atom as a necessary threshold when other zero-loss atoms are present.

\subsection*{Structure, memory and exact realizations}
The full-history finite-obstruction theorem and its dual characterize common bounded-score synthesis on finite-branching interfaces. The reachable-support potential theorem gives necessary and sufficient cost and score certificates at a fixed classical memory budget. Reusable task identities require a different object, the report-resolved Hermitian span, not the average returned state or every state on its Hilbert support. The finite-depth continuation criterion includes standard Borel reports and a physically realizable sharp depth $d^2-1$ for one initial preparation. Existing Hermitian reachability methods are explicitly credited \cite{Tzeng92,Fu24}; neither their dimension count nor generic quantum control is newly invented here.

The finite reset-instrument theorem derives Bellman continuation, a spectral minimax preparation and equality-based memory certificates from the declared failure branches. It includes any finite family of rank-one detectors in arbitrary finite dimension. Separately, two arbitrary-rank projective detectors have a complete integer memory frontier: for $P,Q\ne0$ and $c=\|PQ\|<1$, the full-history value is $1+[2(1-c)]^{-1}$, the one-live-label value is $1+(1-c)^{-1}$, and the exact two-label value undergoes the transition in Theorem~\ref{thm:two-label-exact}. Three fixed-start live labels attain the unrestricted value. The two-label minimax gap shows why a sharp least-favorable preparation need not survive a nonconvex memory restriction. Principal-angle geometry and first-detection theory have their own established antecedents \cite{Halmos69,Kessler21,DelVecchio26}.

The two stopping interfaces remain distinct. An event-stopped observer can reuse a live label for terminal output because the public final event ends computation. A state-stopped observer has an intrinsic fixed halt map and disjoint live/output types. Their counts are respectively $\max(\ell,t)$ and $\ell+t$. A physical port, continuing phase, retained random seed or stopping clock is not an extra free input. Program data, scratch, numerical precision and physical time are separately charged; a finite-label existence bound does not optimize those resources.

The article also includes non-power, regular, stratified and atomic acquired geometry, singular task matrices, noncommuting post-command realizations, calibrated composition, and a fixed-law compact-exhausted predictive quotient. The complete accompanying mathematical companion contains nonintervening directional sensing and earlier exact stopping and absolute-calibration proofs \cite{Companion}. Independent finite-order measurement, two-stage reset communication, rank-constrained discrimination and spectral/dual developments keep their original source identities. They are not inferred from the present stopping-risk approximation. In particular a common executable predictive quotient for an arbitrary unknown-model family, a general characterization of when the geometric profile is exact, and matching program/precision/workspace/time lower bounds remain separate research questions. The current theorems advance the original Foundations problem without identifying any one special realization with its completion.
